\begin{longtable}{TD}
\caption{Influence domain terms.}
\label{tab:terms-influence-domains}\\

\toprule
Term & Definition \\
\midrule
\endfirsthead

\toprule
Term & Definition \\
\midrule
\endhead

\bottomrule
\endfoot

\termrow{influence domain}{
A semantic role assigned to a model variable in the task behavior model.
Influence domains are not metric domains and not telemetry sources: they group
model elements by role rather than by where a signal was observed.
}

\termrow{task identity domain}{
The influence domain containing identifiers and metadata that describe which
task is being executed, such as task kind, version, queue, tenant, producer,
priority, deadline, and attempt number.
}

\termrow{intrinsic feature domain}{
The influence domain describing work contained in the task instance itself, such
as payload size, record count, object count, partition size, input format,
operation mode, or batch size.
}

\termrow{execution context domain}{
The influence domain describing the input-side environment and state under which
the task attempt executes, such as network condition, dependency health, replica
state, disk pressure, cache warmth, worker pressure, retry state, and tenant
pressure.
}

\termrow{resource demand domain}{
The influence domain describing active resource dimensions required or consumed
by the task attempt, such as CPU time, memory peak, disk I/O, network I/O,
runtime allocations, or accelerator usage. Waiting on dependencies is not
active resource demand by itself, although it can contribute to worker occupancy
or risk.
}

\termrow{worker occupancy domain}{
The influence domain describing bounded execution capacity held by the task
attempt, such as worker-slot time, lease time, connection hold time, lock hold
time, semaphore-permit time, or tenant-concurrency time.
}

\termrow{symptom domain}{
The influence domain describing intermediate observed signals that indicate
abnormal or relevant behavior but are not final outcomes and not cost dimensions
by themselves. Examples include dependency wait spikes, retry bursts, timeout
warnings, queue-delay spikes, and lock-contention signals.
}

\termrow{outcome domain}{
The influence domain describing how the task attempt completed or transitioned
in its lifecycle, such as success, failure, timeout, cancellation, retryable
error, non-retryable error, dead-lettering, or lease expiration.
}

\termrow{reliability diagnostic domain}{
The influence domain describing estimate reliability and model-quality evidence,
such as residual magnitude, missing features, sample scarcity, drift, tail ratio,
multimodality, calibration error, or confidence.
}

\termrow{risk domain}{
The influence domain describing the consequence-weighted danger of
underestimating, misclassifying, or mishandling a task attempt. Examples include
queue starvation, retry storm, SLO violation, replication lag, tenant impact,
dependency overload, and cascading overload.
}

\termrow{profiling-policy trigger domain}{
The influence domain describing variables or decisions that select an
observation strategy, such as minimal accounting, detailed accounting, sampled
tracing, tail-aware profiling, or diagnostic profiling.
}

\termrow{one-domain rule}{
The rule that each Task Behavior Graph node has one primary influence domain. A
single raw observation may map to several model variables when it supports
several roles.
}

\termrow{anti-double-counting rule}{
The rule that a signal must be assigned to its model role before it is
aggregated, scored, or used as evidence for another role.
}

\end{longtable}
